\subsection{Conjugación de amplitudes e invariante orientado}
\label{sec:ii09-conjugacion-amplitudes}

La realización compleja conserva los tres ángulos de mezcla y
la fase de interferencia como coordenadas explícitas. Escribamos
\(V(\boldsymbol\vartheta,\delta)
=R_{23}U_{13}(\delta)R_{12}\), con las rotaciones y la fase
definidas en el desarrollo unitario. Los argumentos de las
funciones trigonométricas son sus representantes en radianes.

\begin{proposition}[Conjugación en la realización de mezcla]
\label{prop:ii09-conjugacion-unitaria}
Para los tres ángulos reales y la fase real declarados,
\[
 V(\boldsymbol\vartheta,-\delta)
       =\overline{V(\boldsymbol\vartheta,\delta)},
 \qquad
 J(\boldsymbol\vartheta,-\delta)
       =-J(\boldsymbol\vartheta,\delta).
\]
El invariante \(J\) conserva además su valor bajo los cambios
diagonales de fase de las bases cargadas.
\end{proposition}
\begin{proof}
Las matrices \(R_{12}\) y \(R_{23}\) tienen entradas reales.
En \(U_{13}\), el cambio \(\delta\mapsto-\delta\) conjuga
los factores exponenciales. La conjugación del producto da la
primera identidad. En el cuarteto
\(V_{11}V_{22}\overline{V_{12}}\overline{V_{21}}\), los
factores diagonales de cada fila y columna se cancelan. La
conjugación invierte el signo de su parte imaginaria. Se obtiene
equivalentemente el resultado sustituyendo en
\[
 J=c_{12}c_{23}c_{13}^{2}s_{12}s_{23}s_{13}\sin\delta.
\]
\end{proof}

En el sector de tres generaciones, el valor interno no nulo de
\(J\) expresa la obstrucción a realizar todas las amplitudes
mediante una matriz real por cambios diagonales de fase. Esta
es la lectura de violación de carga--paridad de la matriz
construida. La lectura de rutas conserva, a su vez, la operación
discreta de conjugación y paridad y su representación. Su
compatibilidad completa se expresa mediante el cuadrado
\[
 \mathcal V(CP_{\rm HMT}\rho)
  =D_u\,\overline{\mathcal V(\rho)}\,D_d^{\dagger},
\]
con los espacios de representación y los cambios de fase
declarados. La identidad de conjugación probada aquí fija
exactamente la acción que debe transportar ese cuadrado.
